% !TEX program = xelatex
\documentclass[a4paper,10pt]{oblivoir}
\setmainhangulfont{Noto Serif CJK KR}
\setsanshangulfont{Noto Sans CJK KR}
\usepackage{lipsum}
\begin{document}
\title{학위 논문 형식 예제}\author{홍길동}\date{2024년 1월}
\maketitle
\tableofcontents
\chapter{서론}
이 문서는 oblivoir 문서 클래스를 사용한 한국어 보고서 예제입니다. 국문 조판에서 흔히 쓰이는 장, 절, 각주\footnote{각주 예시입니다.}를 포함합니다.
\section{연구 배경}
한국어 문단은 글자 사이에서 줄바꿈이 일어날 수 있으며, 문장 부호의 처리 방식이 엔진마다 다를 수 있습니다. 이 문단은 여러 줄에 걸쳐 조판되도록 충분히 길게 작성되었습니다. 다시 한 번 같은 내용을 반복하여 줄바꿈을 유도합니다. 한국어 문단은 글자 사이에서 줄바꿈이 일어날 수 있습니다.
\chapter{결론}
본 예제는 \LaTeX{} 엔진 간 차이를 측정하기 위한 것입니다.
\end{document}
